% ---
% id: receiver/def:fock
% kind: def
% title: "Fock space and the second-quantization functor"
% summary: >-
%   The Fock space over a finite-dimensional space is the direct sum of its
%   symmetric photon-number sectors, and a linear map lifts to it sector by
%   sector. Every photon-number block of the receiver's measurement is
%   assembled from blocks of such lifts.
% prominence: standard
% status: accepted
% uses: []
% mentions: [lem:fock-functor, def:physmodel]
% assumes: []
% refs: []
% provenance: newly-derived
% lean:
% ---

The lift is written for a map between two spaces as well as for an operator on
one; the two-space form is used for coordinate inclusions of mode spaces,
among them the inclusion of the \tref{declared input modes} into the dilated
mode space (Definition~\ref{def:physmodel}).

\begin{definition}
\label{def:fock}
For a finite-dimensional Hilbert space $\mathcal{K}_0$ the \term[The direct sum over n of the symmetric n-photon subspaces of a finite-dimensional mode space.]{Fock space} is
$\mathcal{F}(\mathcal{K}_0)=\bigoplus_{n\ge0}\Symn{n}(\mathcal{K}_0)$, where
$\Symn{n}(\mathcal{K}_0)\subseteq\mathcal{K}_0^{\otimes n}$ is the subspace
invariant under every permutation of the $n$ tensor factors and
$\Symn{0}(\mathcal{K}_0)=\mathbb{C}\ket{\vac}$. For a linear operator $X$ on
$\mathcal{K}_0$ put
\begin{equation}\label{eq:fockfunctor}
  \Fock(X)\coloneqq\bigoplus_{n\ge0}\big(X^{\otimes n}\big)\big|_{\Symn{n}},
  \qquad \Fock(X)\big|_{\Symn{0}}=1 .
\end{equation}
This is well defined because $X^{\otimes n}$ commutes with every permutation of
the factors and therefore preserves $\Symn{n}(\mathcal{K}_0)$
(Lemma~\ref{lem:fock-functor}). The same formula defines
$\Fock(X):\mathcal{F}(\mathcal{K}_0)\to\mathcal{F}(\mathcal{K}_1)$ for a linear
map $X:\mathcal{K}_0\to\mathcal{K}_1$ between two spaces, because
$X^{\otimes n}$ intertwines the permutation actions and therefore carries
$\Symn{n}(\mathcal{K}_0)$ into $\Symn{n}(\mathcal{K}_1)$. Occupation vectors
$\vec m\in\mathbb{N}_0^{\dim\mathcal{K}_0}$ with $|\vec m|=n$ index the
orthonormal basis $\{\ket{\vec m}\}$ of $\Symn{n}(\mathcal{K}_0)$ obtained by
symmetrizing and normalizing $\ket{e_1}^{\otimes m_1}\otimes\cdots$ in a fixed
orthonormal basis $\{\ket{e_j}\}$ of $\mathcal{K}_0$; $a_j$ and $a_j^\dagger$
denote the annihilation and creation operators of $\ket{e_j}$.
\end{definition}
